% GENERATED FILE. Edit canonical lessons, then run npm run generate:books.
% canonical-source-hash: fnv1a64:c7f71549c4ca1e6d
% canonical-lessons: ML-C558-kusalam, ML-C558-balyam, ML-C558-koncuka, ML-C558-mutirnna, ML-C558-irattakkutti

\chapter{Youth and Old Age}
\label{ch:ml-youth-and-old-age}
\hlchaptermodality{558}

\begin{chapteropening}
\textbf{By the end of this chapter:} \emph{I can say pleasantries, early childhood, to prattle, grown-up and a twin.}

Complete the last lesson of chapter 558: I can say pleasantries, early childhood, to prattle, grown-up and a twin.
\end{chapteropening}

\section[\texorpdfstring{\ml{കുശലം} (kuśalaṁ)}{കുശലം (kuśalaṁ)}]{\ml{കുശലം} (kuśalaṁ) — pleasantries; small talk}

\label{lesson:ML-C558-kusalam}



\begin{quote}
Before the new one: say the Malayalam for to cuddle, then the Malayalam for a get-together.
\end{quote}

\subsection*{You'll want to know: \ml{കുശലം}}
\textbf{\ml{കുശലം}} — \emph{kuśalaṁ} — \textquotedblleft{}pleasantries; small talk\textquotedblright{}.

\begin{quote}
\ml{വഴിയിൽ} \ml{കണ്ട} \ml{അയൽക്കാരനോട്} \ml{അച്ഛൻ} \ml{കുശലം} \ml{ചോദിച്ചു}.
\end{quote}

\begin{quote}
\emph{vaḻiyil kaṇṭa ayalkkāranōṭŭ acchan kuśalaṁ cōdiccu.}
\end{quote}

\begin{quote}
Father asked after the neighbour he met on the road.
\end{quote}

\ml{കുശലം} \ml{ചോദിക്കുക}, to ask after someone, is the set phrase.

\begin{cousinweb}
\textbf{\ml{കുശലം}} (\emph{kuśalaṁ}) is borrowed from Sanskrit \emph{kuśala}, \textquotedblleft{}right, proper, healthy\textquotedblright{} (CDIAL 3365): to ask the \ml{കുശലം} is to ask whether all is well.
\end{cousinweb}

\begin{grammarlens}[title={What you've built}]
One more word for describing what you lived and felt.
\end{grammarlens}

\subsection*{Your turn}
\begin{itemize}
\raggedright
  \item \emph{Say it:} \emph{kuśalaṁ}
  \item \emph{Say it:} \emph{kuśalaṁ}, once more
  \item \emph{Say it:} \emph{vaḻiyil kaṇṭa ayalkkāranōṭŭ acchan kuśalaṁ cōdiccu.} — Father asked after the neighbour he met on the road.
  \item \emph{Recall:} say the Malayalam for to cuddle, then the Malayalam for a get-together, then say \emph{kuśalaṁ} again
\end{itemize}

\begin{tcolorbox}[breakable,colback=teal!4,colframe=teal!35!black,title={Before you move on}]
What does \ml{കുശലം} mean? (\textquotedblleft{}Pleasantries; small talk\textquotedblright{}.) Say it once more.
\end{tcolorbox}
\section[\texorpdfstring{\ml{ബാല്യം} (bālyaṁ)}{ബാല്യം (bālyaṁ)}]{\ml{ബാല്യം} (bālyaṁ) — early childhood}

\label{lesson:ML-C558-balyam}



\begin{quote}
Before the new one: say the Malayalam for a get-together, then the Malayalam for pleasantries.
\end{quote}

\subsection*{You'll want to know: \ml{ബാല്യം}}
\textbf{\ml{ബാല്യം}} — \emph{bālyaṁ} — \textquotedblleft{}early childhood\textquotedblright{}.

\begin{quote}
\ml{എന്റെ} \ml{ബാല്യം} \ml{ഗ്രാമത്തിലായിരുന്നു}.
\end{quote}

\begin{quote}
\emph{enṟe bālyaṁ grāmattilāyirunnu.}
\end{quote}

\begin{quote}
My early childhood was spent in the village.
\end{quote}

A word of writing; in speech one says \ml{കുട്ടിക്കാലം}.

\begin{cousinweb}
\textbf{\ml{ബാല്യം}} (\emph{bālyaṁ}) is borrowed from Sanskrit \emph{bālya}, \textquotedblleft{}boyhood, childhood\textquotedblright{} (Monier-Williams), the noun of \emph{bāla}, \textquotedblleft{}a child\textquotedblright{}.
\end{cousinweb}

\begin{grammarlens}[title={What you've built}]
One more word for describing what you lived and felt.
\end{grammarlens}

\subsection*{Your turn}
\begin{itemize}
\raggedright
  \item \emph{Say it:} \emph{bālyaṁ}
  \item \emph{Say it:} \emph{bālyaṁ}, once more
  \item \emph{Say it:} \emph{enṟe bālyaṁ grāmattilāyirunnu.} — My early childhood was spent in the village.
  \item \emph{Recall:} say the Malayalam for a get-together, then the Malayalam for pleasantries, then say \emph{bālyaṁ} again
\end{itemize}

\begin{tcolorbox}[breakable,colback=teal!4,colframe=teal!35!black,title={Before you move on}]
What does \ml{ബാല്യം} mean? (\textquotedblleft{}Early childhood\textquotedblright{}.) Say it once more.
\end{tcolorbox}
\section[\texorpdfstring{\ml{കൊഞ്ചുക} (koñcuka)}{കൊഞ്ചുക (koñcuka)}]{\ml{കൊഞ്ചുക} (koñcuka) — to prattle; to talk baby talk}

\label{lesson:ML-C558-koncuka}



\begin{quote}
Before the new one: say the Malayalam for pleasantries, then the Malayalam for early childhood.
\end{quote}

\subsection*{You'll want to know: \ml{കൊഞ്ചുക}}
\textbf{\ml{കൊഞ്ചുക}} — \emph{koñcuka} — \textquotedblleft{}to prattle; to talk baby talk\textquotedblright{}.

\begin{quote}
\ml{കുഞ്ഞ്} \ml{അമ്മയോട്} \ml{കൊഞ്ചി}.
\end{quote}

\begin{quote}
\emph{kuññŭ ammayōṭŭ koñci.}
\end{quote}

\begin{quote}
The baby babbled sweetly to its mother.
\end{quote}

Past: \ml{കൊഞ്ചി}, \emph{koñci}.

Also the cooing of grown-ups to a baby; \ml{കൊഞ്ചൽ} is the noun.

\begin{cousinweb}
\textbf{\ml{കൊഞ്ചുക}} (\emph{koñcuka}) is native Dravidian: DEDR 2043 gives the Malayalam verb as to prattle and to fondle, as with a child (DEDR 2043).
\end{cousinweb}

\begin{grammarlens}[title={What you've built}]
One more word for describing what you lived and felt.
\end{grammarlens}

\subsection*{Your turn}
\begin{itemize}
\raggedright
  \item \emph{Say it:} \emph{koñcuka}
  \item \emph{Say it:} \emph{koñcuka}, once more
  \item \emph{Say it:} \emph{kuññŭ ammayōṭŭ koñci.} — The baby babbled sweetly to its mother.
  \item \emph{Recall:} say the Malayalam for pleasantries, then the Malayalam for early childhood, then say \emph{koñcuka} again
\end{itemize}

\begin{tcolorbox}[breakable,colback=teal!4,colframe=teal!35!black,title={Before you move on}]
What does \ml{കൊഞ്ചുക} mean? (\textquotedblleft{}To prattle; to talk baby talk\textquotedblright{}.) Say it once more.
\end{tcolorbox}
\section[\texorpdfstring{\ml{മുതിർന്ന} (mutirnna)}{മുതിർന്ന (mutirnna)}]{\ml{മുതിർന്ന} (mutirnna) — grown-up; senior}

\label{lesson:ML-C558-mutirnna}



\begin{quote}
Before the new one: say the Malayalam for early childhood, then the Malayalam for to prattle.
\end{quote}

\subsection*{You'll want to know: \ml{മുതിർന്ന}}
\textbf{\ml{മുതിർന്ന}} — \emph{mutirnna} — \textquotedblleft{}grown-up; senior\textquotedblright{}.

\begin{quote}
\ml{വീട്ടിലെ} \ml{മുതിർന്ന} \ml{ആളുകളെ} \ml{കുട്ടികൾ} \ml{ബഹുമാനിക്കണം}.
\end{quote}

\begin{quote}
\emph{vīṭṭile mutirnna āḷukaḷe kuṭṭikaḷ bahumānikkaṇaṁ.}
\end{quote}

\begin{quote}
Children should respect the grown-ups at home.
\end{quote}

\ml{മുതിർന്നവർ} means the elders, the adults.

\begin{cousinweb}
\textbf{\ml{മുതിർന്ന}} (\emph{mutirnna}) is native Dravidian, the adjective of \ml{മുതിരുക}, \textquotedblleft{}to grow ripe, to mature\textquotedblright{}, which DEDR 4954 lists with \ml{മുതു}, \textquotedblleft{}old\textquotedblright{} (DEDR 4954).
\end{cousinweb}

\begin{grammarlens}[title={What you've built}]
One more word for describing what you lived and felt.
\end{grammarlens}

\subsection*{Your turn}
\begin{itemize}
\raggedright
  \item \emph{Say it:} \emph{mutirnna}
  \item \emph{Say it:} \emph{mutirnna}, once more
  \item \emph{Say it:} \emph{vīṭṭile mutirnna āḷukaḷe kuṭṭikaḷ bahumānikkaṇaṁ.} — Children should respect the grown-ups at home.
  \item \emph{Recall:} say the Malayalam for early childhood, then the Malayalam for to prattle, then say \emph{mutirnna} again
\end{itemize}

\begin{tcolorbox}[breakable,colback=teal!4,colframe=teal!35!black,title={Before you move on}]
What does \ml{മുതിർന്ന} mean? (\textquotedblleft{}Grown-up; senior\textquotedblright{}.) Say it once more.
\end{tcolorbox}
\section[\texorpdfstring{\ml{ഇരട്ടക്കുട്ടി} (iraṭṭakkuṭṭi)}{ഇരട്ടക്കുട്ടി (iraṭṭakkuṭṭi)}]{\ml{ഇരട്ടക്കുട്ടി} (iraṭṭakkuṭṭi) — a twin}

\label{lesson:ML-C558-irattakkutti}



\begin{quote}
Before the new one: say the Malayalam for to prattle, then the Malayalam for grown-up.
\end{quote}

\subsection*{You'll want to know: \ml{ഇരട്ടക്കുട്ടി}}
\textbf{\ml{ഇരട്ടക്കുട്ടി}} — \emph{iraṭṭakkuṭṭi} — \textquotedblleft{}a twin\textquotedblright{}.

\begin{quote}
\ml{എന്റെ} \ml{സുഹൃത്തിന്} \ml{ഇരട്ടക്കുട്ടികൾ} \ml{ജനിച്ചു}.
\end{quote}

\begin{quote}
\emph{enṟe suhṛttinŭ iraṭṭakkuṭṭikaḷ janiccu.}
\end{quote}

\begin{quote}
My friend had twins.
\end{quote}

Plural: \ml{ഇരട്ടക്കുട്ടികൾ}, \emph{iraṭṭakkuṭṭikaḷ}.

Usually plural: \ml{ഇരട്ടക്കുട്ടികൾ}, twins. \ml{ഇരട്ട}, double, is already familiar.

\begin{cousinweb}
\textbf{\ml{ഇരട്ടക്കുട്ടി}} (\emph{iraṭṭakkuṭṭi}) is native Dravidian, a compound of \ml{ഇരട്ട}, \textquotedblleft{}double\textquotedblright{} (DEDR 474), and \ml{കുട്ടി}, \textquotedblleft{}a child\textquotedblright{} (DEDR 1670).
\end{cousinweb}

\begin{grammarlens}[title={What you've built}]
One more word for describing what you lived and felt.
\end{grammarlens}

\subsection*{Your turn}
\begin{itemize}
\raggedright
  \item \emph{Say it:} \emph{iraṭṭakkuṭṭi}
  \item \emph{Say it:} \emph{iraṭṭakkuṭṭi}, once more
  \item \emph{Say it:} \emph{enṟe suhṛttinŭ iraṭṭakkuṭṭikaḷ janiccu.} — My friend had twins.
  \item \emph{Recall:} say the Malayalam for to prattle, then the Malayalam for grown-up, then say \emph{iraṭṭakkuṭṭi} again
\end{itemize}

\begin{tcolorbox}[breakable,colback=teal!4,colframe=teal!35!black,title={Before you move on}]
What does \ml{ഇരട്ടക്കുട്ടി} mean? (\textquotedblleft{}A twin\textquotedblright{}.) Say it once more.
\end{tcolorbox}
